% ═════════════════════════════════════════════════════════════════════════════
\section{Abstraction and Network Invariants}
\label{sec:abstraction}
% ═════════════════════════════════════════════════════════════════════════════

Abstraction-based methods address the parameterized problem by reducing it to a
finite, non-parameterized one. The idea is to construct an abstract system,
typically a small, fixed-size system, that \emph{over-approximates} the
behavior of the full parameterized family. If a safety property holds in the
abstract system, it holds in all concrete instances~\cite{WL89,PXZ02,ZP04}.
The methods in this section differ in what the abstract system describes: the
environment of a single process, the number of processes in each local state,
or the relationship between a reference process and the others.

\subsection{Network Invariants}

A particularly influential instantiation of this idea is the \emph{network
invariant} method~\cite{WL89,KM89}, further developed
in~\cite{KPSZ02,LHR97}. Consider a system consisting of $n$ copies of a
process template running in parallel. A network invariant is a single
process~$I$ (or small composition of processes) that captures the collective
behavior of all $n-1$ processes that a given process interacts with.
Verifying the property then reduces to checking the two-process composition
of one concrete process with the invariant~$I$, regardless of~$n$. If this
finite system satisfies the property, and~$I$ is a valid over-approximation of
the environment, then all $n$-process instances satisfy it too.

Validity of the invariant is established by induction, as described in
Section~\ref{sec:early}: one checks that $I$ over-approximates a single copy of
the template, and that the composition of~$I$ with one more copy is again
over-approximated by~$I$. Both checks are finite-state problems. The notion of
over-approximation is a preorder on processes, such as trace inclusion or
simulation, that is preserved by parallel composition and that preserves the
properties of interest.

As an illustration, consider the running example from the point of view of
one process. The other processes matter to it only through the question of
whether one of them is in~$\Crit$. A candidate invariant~$I$ has two states,
\emph{free} and \emph{busy}. It can move from free to busy only when the
process it is composed with is not in~$\Crit$, and it can move from busy to
free at any time. The composition of one process with~$I$ satisfies mutual
exclusion, since the process enters~$\Crit$ only when $I$ is free, and $I$
becomes busy only when the process is not in~$\Crit$. The composition of~$I$
with one more process, observed through the same question, behaves again
like~$I$, which establishes validity.

Finding the right invariant automatically is the central challenge. Approaches
range from manual construction guided by domain knowledge, to automated
inference using abstract interpretation, to counterexample-guided abstraction
refinement (CEGAR) adapted to the parameterized setting. In the CEGAR approach,
one begins with a coarse invariant, checks the finite abstract system, and if
a spurious counterexample is found, refines the invariant to rule it out.
Lesens, Halbwachs, and Raymond~\cite{LHR97} compute invariants for linear
networks by iterating the composition and using an approximation that
enforces convergence. Kesten, Pnueli, Shahar, and Zuck~\cite{KPSZ02} developed
a deductive framework for network invariants, based on abstraction relations
between systems, and applied it to a range of examples.

\subsection{Counter Abstraction}

The work of~\cite{PXZ02} combines abstraction with
\emph{$(0,1,\infty)$-counter abstraction}, which tracks for each local state
whether zero, one, or many processes occupy it. Since the number of local
states is finite, the abstract system is finite. Its transitions are
over-approximations of the concrete ones: for instance, when a process leaves
a local state whose counter is~$\infty$, the counter may become either~$1$
or remain~$\infty$.

In the running example, the abstract state is a triple of counter values for
$\Idle$, $\Wait$, and~$\Crit$, and the bad states are those in which the
counter for~$\Crit$ is~$\infty$. A process can enter~$\Crit$ only if the
counter for~$\Crit$ is~$0$, which sets it to~$1$, and leaving~$\Crit$ sets it
back to~$0$. The counter for~$\Crit$ therefore never reaches~$\infty$, and
mutual exclusion holds in the abstract system, which has at most 27 states.

This finite abstraction is often coarse but captures enough structure to
verify mutual exclusion and other coordination properties automatically, and
the abstraction can be systematically refined when needed. The main
contribution of~\cite{PXZ02} concerns liveness. The abstraction loses the
identity of processes, so a fairness requirement for each individual process
cannot be expressed directly. The method adds to the abstract system
fairness conditions derived from the concrete system, which are strong enough
to prove properties such as the absence of starvation.

\subsection{Environment Abstraction}

\emph{Environment abstraction}~\cite{CTV06,CTV08} describes the system from
the point of view of a single reference process. An abstract state consists of
the local state of the reference process together with the truth values of a
fixed set of predicates that describe the other processes in relation to it,
for example whether some process with a smaller index is in its critical
section, or whether some other process holds a message addressed to the
reference process. Because all processes are identical, it suffices to
consider one reference process, and the abstract transition relation can be
computed once from the template. The
approach has been applied to mutual exclusion algorithms, such as the Bakery
and Szymanski protocols, and to cache coherence protocols.

\subsection{Compositional Reasoning for Protocols}

Cache coherence protocols have been a main application of abstraction
methods. McMillan~\cite{McM01} verified the FLASH protocol for an arbitrary
number of processors by compositional model checking, combining symmetry,
case splitting on the involved processors, and abstraction of the remaining
ones, with lemmas supplied by the user. Chou, Mannava, and Park~\cite{CMP04}
turned this approach into a simple and widely used method, often called the
\emph{CMP method}. A few processes are kept concrete, and all others are
replaced by a single abstract process that can perform any action. This
abstraction is usually too coarse, and spurious counterexamples are
eliminated by adding noninterference lemmas, which constrain the actions of
the abstract process and are themselves verified in the abstract model.
Talupur and Tuttle~\cite{TT08} showed that many of the required lemmas can be
derived from the \emph{message flows} of a protocol, the sequences of messages
exchanged to complete a transaction, which designers typically document.

Predicate abstraction has also been lifted to the parameterized setting.
\emph{Indexed predicates} contain free variables that range over process
indices, and the abstraction computes invariants that are universally
quantified over these variables~\cite{LB07}. This makes it possible to infer
properties such as ``for every process~$i$, if $i$ is in its critical section
then the shared lock is held by~$i$''.

\subsection{Thread-Modular Reasoning}

For concurrent programs with an unbounded number of identical threads,
\emph{thread-modular} reasoning verifies a single thread against an
assumption on the effect of the other threads, and checks that each thread
respects this assumption~\cite{FQ03}. Since all threads are identical, a
single assumption serves for all of them, which makes the method naturally
parameterized. The precision of thread-modular reasoning can be increased by
reasoning about~$k$ threads at a time instead of one, which yields a hierarchy
of increasingly precise proof methods~\cite{HMP17}. \emph{Proof
spaces}~\cite{FKP15} provide a related framework, in which correctness proofs
for an unbounded number of threads are built from a finite set of Hoare
triples that are combined using rules that exploit the symmetry of the
threads.

\subsection{Relation to View Abstraction}

The network invariant approach is closely related to view abstraction
(Section~\ref{sec:cutoff}) in spirit: both methods replace the unbounded
environment of a process by a finite abstract summary. The key difference is
in how that summary is constructed and updated. A network invariant is a
single process that must be supplied, or computed and refined, as a separate
object, and its validity is checked by induction; view abstraction
constructs the summary on the fly from the concrete reachability analysis,
refining it incrementally as the precision parameter~$k$ increases. Counter
abstraction and environment abstraction lie in between: the abstract domain is
fixed in advance, but the abstract system is computed automatically from the
template.
